% =====================================================================
% 第三部分 · 习近平新时代中国特色社会主义思想
% =====================================================================

\zzchapter{第三部分\quad 习近平新时代中国特色社会主义思想}

\zzsection{一、法}

\zzsub{全面依法治国}

\zzcard{全面依法治国}{%
  \kw{“李倩只管四种手机”}\\
  \bk{科学立法}——\kw{前提}；\bk{严格执法}——\kw{关键}；\bk{司法保障}——\kw{重点}；\bk{全民守法}——\kw{基础}。
}

\zzcard{权利保障}{%
  \kw{“县里行司，前要管后”}\\
  \bk{宪法保障}——\kw{前提和基础}；\bk{立法保障}——\kw{重要条件}；\bk{行政保护}——\kw{关键环节}；\bk{司法保障}——\kw{最后防线}。
}

\zzcard{法律的运行}{%
  \kw{“定期、只要、用活、手机”}\\
  \bk{法律制定}——\kw{起始性和关键性环节}；\bk{法律执行}——\kw{重要环节}；\bk{法律适用}——\kw{专门活动}；\bk{法律遵守}——\kw{基本途径}。
}

\zzcard{我国公民的最高法律义务}{%
  我国公民的\kw{最高法律义务}是：\bk{维护国家统一和全国各民族团结}。
}

\zzsub{建设中特法治体系}

\zzcard{建设中特法治体系}{%
  \kw{“高中前辈用医保买了药膳和煎包”}\\
  \bk{高效的法治实施体系}——\kw{重点}；\bk{完备的法律规范体系}——\kw{前提}；\bk{有力的法治保障体系}——\kw{重要依托}；\bk{完善的党内法规体系}——\kw{本质要求}；\bk{严密的法治监督体系}——\kw{重要保障}。
}

\zzcard{法律至上}{%
  \kw{“事必为”}：法律至上具体表现为法律的\bk{普遍适用性}、\bk{优先适用性}、\bk{不可违反性}。\\
  \zzwarn{辨析——\bk{普遍适用性}：所有人都必须遵守；\bk{优先适用性}：多种规范矛盾时，优先考虑法律规范的适用；\bk{不可违反性}：必须遵守，违反要受到惩罚。}
}

\zzcard{推进全面依法治国}{%
  \bk{根本目的}——\kw{依法保障人民权益}；\\
  \bk{根本保证}——\kw{坚持党对全面依法治国的领导}；\\
  \bk{总目标}——\kw{建设中国特色社会主义法治体系，建设社会主义法治国家}；\\
  \bk{总抓手}——\kw{中国特色社会主义法治体系}。\\
  \zzsub{五项基本原则}
  ①坚持中国共产党的领导；②坚持人民主体地位；③坚持法律面前人人平等；④坚持依法治国和以德治国相结合；⑤坚持从中国实际出发。\\
  \zzkey{中国特色社会主义的\bk{本质要求和重要保障}：\kw{全面依法治国}。}
}

\zzsection{二、党}

\zzsub{实现党的全面领导}

\zzcard{实现党的全面领导}{%
  \bk{党中央对重大工作的领导体制}，是加强党的全面领导的\kw{制度安排}，是\kw{实现党的全面领导的直接体现}。加强党的全面领导，必须加强党中央对党和国家重大工作的领导。
}

\zzcard{党的纪律}{%
  党的纪律是多方面的，其中\bk{政治纪律}是\kw{最重要、最根本、最关键的纪律}，遵守党的政治纪律是遵守党的全部纪律的\kw{重要基础}。
}

\zzcard{党内监督、制度建设、纪律建设}{%
  ①\bk{党内监督}是\kw{永葆党的肌体健康的生命之源}，在党和国家各种监督形式中是最根本的、第一位的；\\
  ②\bk{制度建设}是贯穿在政治建设、思想建设、组织建设、作风建设、纪律建设之中的；\\
  ③\bk{加强纪律建设}是全面从严治党的\kw{治本之策}。
}

\zzcard{马克思主义科学理论}{%
  党坚定信仰信念、把握历史主动的\kw{根本所在}；党鲜明的\bk{政治品格}和强大的\bk{政治优势}：\kw{拥有马克思主义科学理论指导}。
}

\zzsub{统一战线}

\zzcard{统一战线}{%
  \kw{“党的领导，鲜明本心”}\\
  统一战线是\bk{党领导}的统一战线：\bk{党的领导是统一战线最鲜明的特征}；\bk{坚持党的领导是统一战线最核心最根本的问题}；加强新时代统一战线工作，根本在于坚持党的领导。\\
  党对统一战线的领导主要是\bk{政治领导}，即政治原则、政治方向、重大方针政策的领导。主要体现为：\kw{党委领导}（部门领导×）、\kw{集体领导}（个人领导×）。
}

\zzsub{我们党区别于其他政党}

\zzcard{我们党区别于其他政党}{%
  \begin{zztable2}
    表述 & 内容 \\
    \midrule
    区别于其他任何政党的\bk{显著标志} & \kw{三大作风}（七大） \\
    区别于其他政党的\bk{显著标志} & \kw{勇于自我革命} \\
    区别于其他一切政党\bk{根本的标志} & \kw{全心全意为人民服务} \\
    区别于其他政党的\bk{显著标志}（新大纲） & \kw{人民立场} \\
  \end{zztable2}
  \zzwarn{\bk{人民立场}是中国共产党的\kw{根本政治立场}，是我们党区别于其他政党的\kw{显著标志}。}
}

\zzsub{自我革命}

\zzlongcard{自我革命}{%
  ①中国共产党区别于其他政党的\bk{显著标志}：\kw{勇于自我革命}；\\
  ②党\bk{最鲜明}的品格：\kw{勇于自我革命}；\\
  ③党如何跳出治乱兴衰的历史周期率给出的\bk{第二个答案}：\kw{勇于自我革命}；\\
  ④党\bk{最大的优势}：\kw{勇于自我革命}。\\
  \zzsub{实现伟大自我革命的重要途径}
  \kw{“革净高完”}：\bk{自我净化}、\bk{自我完善}、\bk{自我革新}、\bk{自我提高}。\\
  \zzsub{深入推进党的自我革命需要把握好的九个问题}
  以\bk{坚持党中央集中统一领导}为\kw{根本保证}；以\bk{引领伟大社会革命}为\kw{根本目的}（本伟）；以\bk{习近平新时代中国特色社会主义思想}为\kw{根本遵循}；以\bk{跳出历史周期率}为\kw{战略目标}（战士）；以\bk{解决大党独有难题}为\kw{主攻方向}（攻读）；以\bk{健全全面从严治党体系}为\kw{有效途径}（体校）；以\bk{锻造坚强组织、建设过硬队伍}为\kw{重要着力点}（早点）；以\bk{正风肃纪反腐}为\kw{重要抓手}（抓鸡）；以\bk{自我监督和人民监督相结合}为\kw{强大动力}（自强）。
}

\zzsub{党的领导}

\zzsub{党的领导（万能答案）}
\begin{zzltable2}
    考点 & 结论 \\
    \midrule
    中国特色社会主义\bk{最本质的特征} & \kw{党的领导} \\
    中国特色社会主义制度的\bk{最大优势} & \kw{党的领导} \\
    中国特色社会主义制度体系的\bk{“根”和“源”} & \kw{党的领导}制度 \\
    推进全面依法治国的\bk{根本保证} & \kw{党的领导} \\
    人民当家作主的\bk{根本保证} & \kw{党的领导} \\
    我国社会主义\bk{法治之魂} & \kw{党的领导} \\
    社会主义法治\bk{最根本的保证} & \kw{党的领导} \\
    维护国家安全和社会安定的\bk{根本保证} & 坚持\kw{党}对国家安全工作的\bk{绝对领导} \\
    中华民族伟大复兴的\bk{根本保证} & \kw{党的领导} \\
    我们的法治同西方资本主义国家法治的\bk{最大区别} & \kw{党的领导} \\
    生态文明建设的\bk{根本保证} & 加强\kw{党}对生态文明建设的\bk{全面领导} \\
    进一步全面深化改革、推进中国式现代化的\bk{根本保证} & \kw{党的领导} \\
  \end{zzltable2}
%
  
  \zzwarn{\bk{党中央集中统一领导}是：①\kw{党的领导的最高原则}；②党保持团结统一和强大战斗力、\kw{不断取得胜利的关键所在}。（根本保证党领导，胜利关键党领导）}

\zzsection{三、发展}

\zzsub{新发展理念}

\zzcard{新发展理念}{%
  \kw{“宫本武藏对绿箭心动，用开塞露内调”}\\
  \begin{zztable2}
    理念 & 定位 \\
    \midrule
    共享 & \kw{本质要求} \\
    绿色 & \kw{必要条件} \\
    创新 & \kw{第一动力} \\
    开放 & \kw{必由之路} \\
    协调 & \kw{内在要求} \\
  \end{zztable2}
  \kw{“红绿灯跟导游”}：新发展理念是\kw{指挥棒、红绿灯}，是引领我国发展全局深刻变革的科学指引。必须完整、准确、全面地理解和把握，要从\bk{根本宗旨}上、\bk{问题导向}上、\bk{忧患意识}上把握新发展理念。
}

\zzcard{科学发展观}{%
  \kw{“法医仁心在房间吃鸡”}\\
  \bk{推动经济社会发展}——\kw{第一要义}；\bk{以人为本}——\kw{核心立场}；\bk{统筹兼顾}——\kw{根本方法}；\bk{全面协调可持续}——\kw{基本要求}。
}

\zzsub{新发展格局}

\zzsub{新发展格局}
\begin{zzltable2}
    项目 & 内容 \\
    \midrule
    基本立足点 & \kw{满足国内需求} \\
    主线 & \kw{深化供给侧结构性改革} \\
    最本质特征 & \kw{实现高水平自立自强} \\
    战略方向 & \kw{坚持供给侧结构性改革} \\
    关键 & \kw{经济循环畅通无阻} \\
    战略支撑 & \kw{科技自立自强} \\
    重要前提和保障 & \kw{牢牢守住安全发展这条底线} \\
  \end{zzltable2}
%
  \kw{“店内神仙本子多，房间灌肠留成绩”}\\

\zzsub{高质量发展}

\zzcard{高质量发展}{%
  \kw{“科比公主爱打农药，是歌姬，没有福目”}\\
  \begin{zztable2}
    项目 & 内容 \\
    \midrule
    必由之路 & 加快实现高水平\kw{科技自立自强} \\
    主线 & 深化\kw{供给侧改革} \\
    必然要求 & 推进\kw{农业现代化} \\
    战略基点 & 加快构建\kw{新发展格局} \\
    最终目的 & \kw{人民幸福安康} \\
    根本遵循 & \kw{习近平经济思想} \\
  \end{zztable2}
  \zzkey{\bk{高质量发展}是全面建设社会主义现代化国家的\kw{首要任务}。新时代我国\bk{经济社会发展的鲜明主题}是\kw{高质量发展}；新时代我国\bk{经济发展的基本特征}：由高速增长阶段转向高质量发展阶段。}\par\addvspace{2pt}
  \zzsub{首要任务辨析}
  ①\bk{解放和发展社会生产力}是改革开放的鲜明特征和首要任务；\\
  ②\bk{建设社会主义法治国家}，必须把\kw{全面贯彻实施宪法}作为首要任务；\\
  ③\bk{构建统筹各领域安全的新安全格局}，必须把\kw{维护政治安全}作为应对各领域安全风险挑战的首要任务；\\
  ④\bk{保证全党服从中央、维护党中央权威和集中统一领导}是\kw{党的政治建设的首要任务}。\\
  \zzwarn{党的二十大报告指出，推动经济社会发展\bk{绿色化、低碳化}是实现高质量发展的\kw{关键环节}。}\par\addvspace{2pt}
  \zzwarn{建设\bk{现代化产业体系}，要坚持把发展经济的\kw{着力点放在实体经济上}。}
}

\zzsection{四、帽子题汇总}

\zzsub{制度与核心}

\zzcard{制度辨析}{%
  \bk{政协制度}是\kw{基本政治制度}；\bk{人大制度}是\kw{根本政治制度}。\\
  \zzwarn{易错点：\bk{政协}是\bk{组织}不是\bk{制度}；\bk{政协制度}是\bk{制度}。}
}

\zzcard{精神核心}{%
  \bk{时代精神的核心}：改革创新；\\
  \bk{民族精神的核心}：爱国主义；\\
  \bk{五四精神的核心}：爱国主义。
}

\zzcard{良好的生态环境}{%
  最公平的\bk{公共产品}；最普惠的\bk{民生福祉}。
}

\zzsub{发展相关帽子}

\zzcard{科技自立自强}{%
  \kw{“柯基要两只”}\\
  科技自立自强是：\bk{国家强盛之基、安全之要}；\bk{促进发展大局的根本支撑}；\bk{国家发展的战略支撑}。
}

\zzcard{平等}{%
  平等是\bk{社会主义法律的基本属性}，是\bk{社会主义法治的基本要求}。
}

\zzcard{金融工作}{%
  \bk{根本宗旨}——\kw{金融服务实体经济}；\\
  \bk{永恒主题}——\kw{防控风险}。
}

\zzcard{根基前提辨析}{%
  \kw{“全腹肌、舌吻强、弄射鸡”}\\
  \bk{国家安全}是民族复兴的\kw{根基}；\\
  \bk{社会稳定}是国家强盛的\kw{前提}；\\
  \bk{农业强国}是社会主义现代化强国的\kw{根基}。
}

\zzcard{改善生活}{%
  \kw{“结合改善生活”}：正确处理改革发展稳定关系的\kw{结合点}——\bk{改善生活}。
}

\zzsub{共建共治共享与改革方法}

\zzcard{共建共治共享}{%
  \kw{“剑姬只管橡木”}：\bk{共建}是\kw{基础}，\bk{共治}是\kw{关键}，\bk{共享}是\kw{目标}。
}

\zzcard{摸石头过河}{%
  \kw{“摸石头改革”}：富有中国特色、符合中国国情的\kw{改革方法}——\bk{摸石头过河}。
}

\zzcard{构建更高水平开放型经济体制}{%
  我国构建更高水平开放型经济体制的\bk{重点}是\kw{制度型开放}。
}

\zzcard{决定社会主义文化事业前途命运的关键}{%
  决定社会主义文化\bk{事业}前途命运的关键是：\kw{坚持为人民服务、为社会主义服务的根本方向}。
}

\zzsub{区域辨析}

\zzcard{区辨析}{%
  \begin{zztable2}
    区域 & 定位 \\
    \midrule
    京津冀 & 中国式现代化建设的\bk{先行区、示范区} \\
    浦东新区 & 社会主义现代化建设\bk{引领区} \\
    浙江 & \bk{高质量发展建设共同富裕示范区} \\
    深圳 & \bk{中国特色社会主义先行示范区} \\
    福建 & \bk{两岸融合发展示范区} \\
  \end{zztable2}
}

\zzsub{经济制度相关}

\zzcard{四个全面战略布局}{%
  \bk{全面建设社会主义现代化国家}（\kw{引领地位}）、\bk{全面深化改革}、\bk{全面依法治国}、\bk{全面从严治党}。
}

\zzcard{民营经济}{%
  是推进中国式现代化的\kw{生力军}（\bk{不等于}主力军）。\\
  民营经济是\bk{非公有制经济}，是\bk{我国社会主义市场经济的重要组成部分}，但\bk{不是我国社会主义经济的重要组成部分}。
}

\zzcard{国有企业}{%
  \bk{中国特色社会主义的重要物质基础和政治基础}；\\
  我们党执政兴国的\kw{重要支柱和依靠力量}。
}

\zzcard{坚定制度自信}{%
  坚定中国特色社会主义制度自信，\kw{首先要坚定对中国特色社会主义政治制度的自信}。
}

\zzcard{一根红线}{%
  \bk{贯穿中国革命道德始终的一根红线}：\kw{全心全意为人民服务}；\\
  \bk{贯穿习近平新时代中国特色社会主义思想的一根红线}：\kw{人民至上}。
}

\zzsub{一国两制与关键一招}

\zzcard{一国两制}{%
  \bk{维护国家主权、安全、发展利益}是：我国\kw{对外工作的出发点和落脚点}、中国外交的\bk{神圣使命}、\kw{一国两制的最高原则}。\\
  \zzwarn{一国两制的\bk{根本宗旨}：维护国家主权、安全、发展利益，保持香港、澳门长期繁荣稳定。}
}

\zzcard{关键一步，关键一招}{%
  \kw{“小步放大招”}\\
  \bk{全面建成小康社会}是迈向中华民族伟大复兴的\kw{关键一步}；\\
  \bk{改革开放}是决定当代中国命运的\kw{关键一招}，也是决定实现“两个一百年”奋斗目标、实现中华民族伟大复兴的\kw{关键一招}。
}

\zzcard{新型城镇化}{%
  深入实施以人为本的新型城镇化战略，要把\kw{推进农业转移人口市民化}作为新型城镇化\kw{首要任务}。
}

\zzcard{中国特色社会主义进入新时代判断依据}{%
  \kw{“基本矛盾”}：中国特色社会主义进入新时代重大判断的\bk{基本依据}是——\kw{我国社会主要矛盾发生新变化}。
}

\zzsub{法治与出发点落脚点}

\zzcard{坚持法治国家、法治政府、法治社会一体建设}{%
  \kw{“国标争点机会”}\\
  \bk{法治国家}是法治建设的\kw{目标}；\\
  \bk{法治政府建设}是全面依法治国的\kw{重点任务和主体工程}；\\
  \bk{法治社会}是构筑\bk{法治国家的基础}；\bk{法治社会建设}是\kw{基础工程}。
}

\zzlongcard{出发点与落脚点}{%
  \kw{“人民”}\\
  ①\bk{现代化建设}的出发点和落脚点：\kw{实现人民对美好生活的向往}；\\
  ②\bk{发展}的出发点和落脚点：\kw{满足人民对美好生活的新期待}；\\
  ③\bk{社会主义改革开放}的出发点和落脚点：\kw{更好实现和维护人民利益}，让老百姓过上好日子；\\
  ④\bk{民主政治建设}的出发点和落脚点：\kw{实现最广大人民根本利益}；\\
  ⑤\bk{文艺和文艺工作}的出发点和落脚点：\kw{满足人民精神文化需求}；\\
  ⑥\bk{全面建设社会主义现代化国家}的出发点和落脚点：\kw{让人民生活越过越好}；\\
  ⑦\bk{人民幸福}是国家富强、民族振兴的题中之义和必然要求，是国家富强、民族振兴的\kw{根本出发点和落脚点}；\\
  ⑧\bk{进一步全面深化改革}的出发点和落脚点：\kw{促进社会公平正义、增进人民福祉}。\\
  \zzsub{补充}
  ★\bk{全军各项建设}的出发点和落脚点：\kw{提高战斗力}；\\
  ★\bk{我国对外工作}的出发点和落脚点：\kw{维护国家主权、安全、发展利益}。
}

\zzsub{必由之路}

\zzsub{必由之路}
\begin{zzltable2}
    内容 & 必由之路 \\
    \midrule
    坚持和发展中国特色社会主义 & \bk{坚持党的全面领导} \\
    坚持和发展中国特色社会主义 & \bk{改革开放} \\
    实现中华民族伟大复兴 & \bk{中国特色社会主义} \\
    中国人民创造历史伟业 & \bk{团结奋斗} \\
    新时代我国发展壮大 & \bk{贯彻新发展理念} \\
    党永葆生机活力、走好新的赶考之路 & \bk{全面从严治党} \\
    国家繁荣发展 & \bk{开放} \\
    推动高质量发展 & \bk{高水平科技自立自强} \\
    我们攀登世界科技高峰 & \bk{自主创新} \\
    全面准确贯彻“一国两制”方针 & \bk{依法治港治澳} \\
  \end{zzltable2}
%

\zzsection{五、大专题}

\zzsub{外交}

\zzcard{外交}{%
  ①\bk{中国外交政策的宗旨}：\kw{维护世界和平、促进共同发展}；\\
  ②\bk{中国特色大国外交牢牢把握的主线}：\kw{服务民族复兴、促进人类进步}；\\
  ③\bk{周边外交}：理念——\kw{亲诚惠容}；方针——\kw{与邻为善、以邻为伴}；\\
  ④\bk{新型国际关系}：秉持\kw{相互尊重、公平正义、合作共赢}原则（\bk{合作共赢是核心}），走出一条\kw{对话而不对抗、结伴而不结盟}的国与国交往新路。
}

\zzcard{两大历史任务}{%
  \bk{争取民族独立、人民解放}和\bk{实现国家富强、人民幸福}。两大历史任务既有区别，也有联系：\\
  ①\bk{区别}在于：两大历史任务的\kw{主题、内容和实现方式}都不一样；\\
  ②\bk{联系}在于：前一个任务为后一个任务\kw{扫清障碍}，为实现中华民族伟大复兴的中国梦创造必要前提；后一个任务是前一个任务的\kw{最终目的和必然要求}。
}

\zzsub{文化}

\zzcard{文化}{%
  ①\bk{中华优秀传统文化}——我们坚定文化自信的\kw{深厚基础}；\\
  \hspace{1em}\bk{革命文化和社会主义先进文化}（党带领人民在伟大斗争中孕育）——我们坚定文化自信的\kw{坚强基石}（革基）；\\
  \hspace{1em}\bk{中国特色社会主义伟大实践}——我们坚定文化自信的\kw{现实基础}；\\
  ②我们区别于其他国家和民族的\bk{根本特征}：\kw{中华民族创造的优秀传统文化}；\\
  ③\bk{文化自信坚定到底}：坚定中国特色社会主义道路自信、理论自信、制度自信，说到底是\kw{要坚定文化自信}。
}

\zzsub{教育}

\zzcard{教育}{%
  ①\bk{教育强国}：\bk{根本保证}——党对教育事业的全面领导；\bk{根本任务}——立德树人；\bk{根本目标}——为党育人、为国育才；\bk{重要使命}——服务中华民族伟大复兴；\bk{核心功能}——支撑引领中国式现代化；\\
  ②\bk{教育的根本问题、建设教育强国的核心课题}：\kw{培养什么人、怎样培养人、为谁培养人}；\\
  ③\bk{发展教育的根本尺度}：\kw{人民满意}；\\
  ④\bk{教育科技人才}是全面建设社会主义现代化国家的\kw{基础性战略性支撑、先导任务}；\bk{建成教育强国}是实现以中国式现代化全面推进强国建设、民族复兴伟业的\kw{坚实基础、战略支撑、先导任务}；\\
  ⑤\bk{必须坚持}：\bk{科技}——\kw{第一生产力}；\bk{人才}——\kw{第一资源}；\bk{创新}——\kw{第一动力}。
}

\zzsub{中国式现代化}

\zzcard{中国式现代化五个特征}{%
  \kw{“先生奔赴要媳妇，球鞋特显大”}\\
  \begin{zztable2}
    定位 & 特征 \\
    \midrule
    \bk{本质特征}（区别于西方现代化的显著标志） & \kw{全体人民共同富裕的现代化}（奔赴） \\
    \bk{显著特征} & \kw{人口规模巨大的现代化}（特显大） \\
    \bk{突出特征} & \kw{走和平发展道路的现代化}（出道） \\
    \bk{鲜明特点} & \kw{人与自然和谐共生的现代化}（先生） \\
    \bk{崇高追求}（社会主义现代化的本质要求） & \kw{物质文明和精神文明相协调的现代化}（球鞋） \\
  \end{zztable2}
}

\zzcard{中国式现代化的帽子总结}{%
  \bk{中国式现代化创造了人类文明的新形态}：提供了一种全新的现代化模式；是对西方现代化理论和实践的重大超越；为广大发展中国家提供了全新选择。\par\addvspace{2pt}
  ①\bk{中国式现代化}是贯穿各个领域、各个方面、各项工作的一条\kw{红线}；\\
  ②\bk{开放}是中国式现代化的\kw{鲜明标识}；\\
  \hspace{1em}\zzwarn{辨析：\bk{改革}是当代中国人民最鲜明的\kw{精神标识}。}\par\addvspace{2pt}
  ③\bk{发展全过程人民民主}是中国式现代化的\kw{本质要求}；\\
  ④\bk{城乡融合发展}是中国式现代化的\kw{必然要求}。
}

\zzsub{爱国主义}

\zzlongcard{爱国主义}{%
  \kw{“统一团结需要力”}\\
  弘扬新时代爱国主义的\bk{重要着力点和落脚点}是：\kw{维护祖国统一和民族团结}。\\
  \bk{维护祖国统一和民族团结}是我国公民的\kw{最高义务}。\\
  \zzwarn{辨析：①\bk{新时代党的民族工作的出发点和落脚点}是\kw{实现中华民族伟大复兴}；}\par\addvspace{2pt}
  \zzwarn{②\bk{新时代加强爱国主义教育的主题}：旧提法——实现中华民族伟大复兴的中国梦；新提法——\kw{全面建成社会主义现代化强国、实现中华民族伟大复兴}；}\par\addvspace{2pt}
  \zzwarn{③\bk{坚持和发展中国特色社会主义的总任务}：\kw{实现社会主义现代化和中华民族伟大复兴}。}\par\addvspace{2pt}
  \bk{促进中华儿女大团结}，是新时代爱国统一战线的\kw{历史责任}。\\
  \zzkey{坚持爱国和爱党、爱社会主义相统一，是当代中国爱国主义精神最重要的体现。\bk{爱国主义的本质}，就是坚持爱国和爱党、爱社会主义\kw{高度统一}。}
}

\zzsub{总体国家安全观}

\zzcard{总体国家安全观}{%
  \kw{“人质在异国爆射，根治鸡精”}\\
  \begin{zztable2}
    定位 & 内容 \\
    \midrule
    宗旨 & \kw{人民安全} \\
    根本 & \kw{政治安全} \\
    基础 & \kw{经济安全} \\
    保障 & \kw{社会安全} \\
    依托 & \kw{国际安全} \\
  \end{zztable2}
  政治安全与人民安全、国家利益至上的有机统一，确保实现党的长期执政、人民安居乐业、国家长治久安。\\
  ①\bk{政治安全}是维护人民安全和国家利益的\kw{根本保证}；\\
  ②\bk{人民安全}居于\kw{中心地位}，国家安全归根到底是\kw{保障人民利益}；\\
  ③\bk{国家利益至上}是实现政治安全和人民安全的\kw{要求和准则}。\\
  \zzwarn{\bk{健全国家安全体系、加强国家安全制度建设}是保障国家安全的\kw{治本之策}，是维护和塑造国家安全的\kw{重要支撑}。}\par\addvspace{2pt}
  \zzwarn{\bk{国土安全是立国之基}。构建统筹各领域安全的新安全格局，必须把\kw{维护政治安全}作为应对各领域安全风险挑战的\kw{首要任务}；\bk{政治安全}是国家安全的\kw{根本}。}
}

\zzsub{进一步全面深化改革}

\zzlongcard{进一步全面深化改革}{%
  \kw{“自主呻吟、集中重口”}\\
  统筹推进各领域各方面改革开放必须：\bk{以全面深化改革总目标为引领}、\bk{以制度建设为主线}、\bk{以经济体制改革为重点（牵引）}、\bk{以重要领域和关键环节改革为突破口}，使各方面改革协同推进、环环相扣、形成合力。\par\addvspace{2pt}
  ①\bk{主题}：\kw{推进中国式现代化}；\\
  ②\bk{总目标}：继续完善和发展中国特色社会主义制度，推进国家治理体系和治理能力现代化；\\
  ③\bk{六条原则}（严格遵循、长期坚持）：\\
  \hspace{1em}1.\ 以\bk{坚持党的全面领导}为\kw{根本保证}【排在首位】；\\
  \hspace{1em}2.\ 以\bk{坚持以人民为中心}为\kw{价值导向}；\\
  \hspace{1em}3.\ 以\bk{坚持守正创新}为\kw{基本原则}；\\
  \hspace{1em}4.\ 以\bk{坚持以制度建设为主线}为\kw{重要要求}；\\
  \hspace{1em}5.\ 以\bk{坚持全面依法治国}为\kw{有效途径}；\\
  \hspace{1em}6.\ 以\bk{坚持系统观念}为\kw{科学方法}。
}

\zzsub{中华文明的突出特性}

\zzcard{中华文明的突出特性}{%
  \begin{zztable2}
    特性 & 内容 \\
    \midrule
    连续性 & 从根本上决定了中华民族\bk{必然走自己的路} \\
    创新性 & ①从根本上决定了中华民族\bk{守正不守旧、尊古不复古}的进取精神；②决定了中华民族\bk{不惧新挑战、勇于接受新事物}的无畏品格 \\
    统一性 & ①决定了中华民族各民族文化\bk{融为一体}、即使遭遇重大挫折也牢固凝聚；②决定了\bk{国土不可分、国家不可乱、民族不可散、文明不可断}的共同信念；③决定了\bk{国家统一永远是中国核心利益的核心}；④决定了一个坚强统一的国家是各族人民的命运所系 \\
    包容性 & ①决定了中华文化\bk{交往交流交融}的历史取向；②决定了中国\bk{各宗教信仰多元并存}的和谐格局；③决定了中华文化对世界文明\bk{兼收并蓄}的开放胸怀 \\
    和平性 & ①决定了中国始终是\bk{世界和平的建设者、全球发展的贡献者、国际秩序的维护者}；②决定了中国不断追求\bk{文明交流互鉴}而不搞文化霸权；③决定了中国\bk{不会把自己的价值观念与政治体制强加于人}；④决定了中国坚持合作、不搞对抗，\bk{决不搞“党同伐异”的小圈子} \\
  \end{zztable2}
  \zzkey{\kw{“一连创死俩，捅死包小三”}}
}

\zzsub{就业}

\zzcard{就业}{%
  ①\bk{就业（三基）}：①\kw{最基本的民生}；②\kw{最大的民生工程、民心工程、根基工程}；③\kw{劳动者赖以生存和发展的基础、共享经济发展成果的基本}；\\
  ②\bk{收入分配}：是\kw{民生之源}，是\kw{改善民生、实现发展成果由人民共享最重要最直接的方式}（直接分手）；\\
  ③\bk{人力资源供需不匹配}是当前我国就业领域面临的\kw{主要矛盾}。
}

\zzsub{2035 和 2050 目标}

\zzcard{2035 和 2050 目标}{%
  《中共中央关于进一步全面深化改革、推进中国式现代化的决定》指出，进一步全面深化改革的\bk{总目标}是：继续完善和发展中国特色社会主义制度，推进国家治理体系和治理能力现代化。\\
  \zzsub{到 2035 年}
  \bk{全面建成高水平社会主义市场经济体制}，中国特色社会主义制度更加完善，\bk{基本实现}国家治理体系和治理能力现代化，\bk{基本实现社会主义现代化}，为到本世纪中叶全面建成社会主义现代化强国奠定坚实基础。\\
  \zzsub{两步走战略}
  ①从 2020 年到 2035 年\kw{基本实现社会主义现代化}；\\
  ②从 2035 年到本世纪中叶\kw{把我国建成富强民主文明和谐美丽的社会主义现代化强国}，到那时\kw{全体人民共同富裕基本实现}。\\
  \zzwarn{2035 一般都是\bk{“模糊说法”}；2050 一般都是\bk{确切性说法}。}
}

\zzcard{构建高水平社会主义市场经济体制}{%
  \bk{关键}是要处理好\kw{政府和市场的关系}：\\
  ①\bk{市场在资源配置中起决定性作用}——把市场机制能有效调节的经济活动交给市场，推动资源配置实现\kw{效益最大化和效率最优化}；\\
  ②\bk{更好发挥政府作用}——在保证市场发挥决定性作用的前提下，\kw{管好那些市场管不了或管不好的事情}；\\
  ③二者是\kw{有机统一}的，\bk{不是相互否定的}，\bk{不能把二者割裂开来、对立起来}；要用好\kw{“看不见的手”和“看得见的手”}，推动\kw{有效市场和有为政府更好结合}。
}

\zzsub{全面推进乡村振兴}

\zzsub{全面推进乡村振兴}
\begin{zzltable2}
    项目 & 内容 \\
    \midrule
    底线任务／前提 & \kw{巩固脱贫攻坚成果} \\
    总目标 & \kw{农业农村现代化} \\
    总方针 & \kw{坚持农业农村优先发展} \\
    制度保障 & \kw{建立健全城乡融合发展体制机制和政策体系} \\
    首要任务 & \kw{确保重要农产品特别是粮食供给} \\
    “三农”工作的底线 & \kw{保障国家粮食安全和不发生规模性返贫} \\
    总抓手 & \kw{推进乡村全面振兴} \\
    重中之重（关键） & \kw{产业振兴} \\
    头等大事 & \kw{保障粮食和重要农产品稳定安全供给} \\
  \end{zzltable2}
%
  
  \zzwarn{巩固拓展脱贫攻坚成果，是全面推进乡村振兴的\bk{底线任务}。要坚持把\kw{增加农民收入}作为\bk{三农工作中的中心任务}。}\par\addvspace{2pt}
  \zzwarn{\bk{农业强国}是社会主义现代化强国的\kw{根基}；\bk{最艰巨最繁重的任务仍然在农村}；\bk{最广泛最深厚的基础仍然在农村}。}\par\addvspace{2pt}
  \zzsub{2024 最新提法}
  做好 2024 年及今后一个时期“三农”工作：①以\kw{学习运用“千万工程”经验}为【\bk{引领}】；②以\kw{确保国家粮食安全、确保不发生规模性返贫}为【\bk{底线}】。

\zzsub{六个必须坚持}

\zzcard{六个必须坚持}{%
  习近平新时代中国特色社会主义思想的世界观方法论：\kw{坚持人民至上、自信自立、守正创新、问题导向、系统观念、胸怀天下}。\\
  \zzkey{\kw{“本人思念胸境导实的新鲜精目”}}\par\addvspace{2pt}
  \begin{zztable2}
    内容 & 定位与体现 \\
    \midrule
    坚持人民至上 & \kw{根本价值立场}，体现了\bk{历史唯物主义群众史观} \\
    坚持自信自立 & \kw{内在精神特征}，体现了\bk{客观规律性和主观能动性的有机结合} \\
    坚持守正创新 & \kw{鲜明理论品格}，体现了\bk{变与不变、继承与发展的内在联系} \\
    坚持问题导向 & \kw{重要实践要求}，体现了\bk{矛盾的普遍性和客观性} \\
    坚持系统观念 & \kw{基本思想和工作方法}，体现了\bk{辩证唯物主义普遍联系的原理} \\
    坚持胸怀天下 & \kw{中国共产党人的境界格局}，体现了\bk{马克思主义追求人类进步和解放的崇高理想} \\
  \end{zztable2}
}
